\subsection{Beck's Monadicity Theorem}

\begin{thm}\label{thm:beck}{\textit{Beck's Monadicity Theorem}}
	Let $ (F,U,\eta,\epsilon):\mathbb{X}\longrightarrow \mathbb{A}$ be an adjunction, and $T=(T,\eta, \mu)$ be its associated monad.
	The following conditions are equivalent:
	\begin{enumerate}
		\item the functor $ U: \mathbb{A} \longrightarrow \mathbb{X}$ is monadic;
		\item the functor $U$ preserves coequaliser diagrams of the form:
		% https://q.uiver.app/#q=WzAsMyxbMCwwLCJGVUYoWCkiXSxbMSwwLCJGKFgpIl0sWzIsMCwiTChYLFxceGkpIl0sWzAsMSwiXFxlcHNpbG9uX3tGKFgpfSIsMCx7Im9mZnNldCI6LTF9XSxbMCwxLCJGKFxceGkpIiwyLHsib2Zmc2V0IjoxfV0sWzEsMiwiXFxwaV97KFgsIFxceGkpfSJdXQ==
		\[\begin{tikzcd}
			{FUF(X)} & {F(X)} & {L(X,\xi)}
			\arrow["{\epsilon_{F(X)}}", shift left, from=1-1, to=1-2]
			\arrow["{F(\xi)}"', shift right, from=1-1, to=1-2]
			\arrow["{\pi_{(X, \xi)}}", from=1-2, to=1-3]
		\end{tikzcd}\]
		for every $T$-algebra $(X, \xi)$, and for every object $ A$ in $ \mathbb{A}$, the diagram
		\[\begin{tikzcd}
			{FUFU(A)} & {FU(A)} & {A}
			\arrow["{\epsilon_{FU(A)}}", shift left, from=1-1, to=1-2]
			\arrow["{FU(\epsilon_A)}"', shift right, from=1-1, to=1-2]
			\arrow["{ \epsilon_A}", from=1-2, to=1-3]
		\end{tikzcd}\]
		is a coequaliser;
		\item $U$ preserves coequaliser diagrams of the form:
		\[\begin{tikzcd}
			{FUF(X)} & {F(X)} & {L(X,\xi)}
			\arrow["{\epsilon_{F(X)}}", shift left, from=1-1, to=1-2]
			\arrow["{F(\xi)}"', shift right, from=1-1, to=1-2]
			\arrow["{\pi_{(X, \xi)}}", from=1-2, to=1-3]
		\end{tikzcd}\]
		for every $T$-algebra $(X,\xi)$, and $U$ reflects isos;
		\item every diagram 
		\begin{tikzcd}
			{A} & {B} 
			\arrow["{f}", shift left, from=1-1, to=1-2]
			\arrow["{g}"', shift right, from=1-1, to=1-2]
		\end{tikzcd}
		in $ \mathbb{A}$ - where
		\begin{tikzcd}
			{U(A)} & {U(B)} 
			\arrow["{U(f)}", shift left, from=1-1, to=1-2]
			\arrow["{U(g)}"', shift right, from=1-1, to=1-2]
		\end{tikzcd}
		has an absolute coequaliser -  	
		has a coequaliser preserved by $U$, and $U$ reflects isos.
		
	\end{enumerate}
	
\end{thm}

\begin{proof} Since it has been extensively covered that $(1)\iff (2) \iff(3)$, we will only prove $ (1)\implies (4)$.
	Let $ \mathcal{D}$ be a diagram in $ \mathbb{A}$ such that $ U( \mathcal{D}) =$ 
	\begin{tikzcd}
		{X} & {X'}
		\arrow["u", shift left, from=1-1, to=1-2]
		\arrow["v"', shift right, from=1-1, to=1-2]
	\end{tikzcd}
	with in $ \mathbb{X}$. Suppose $ U( \mathcal{D})$ has an absolute coequaliser $q$.
	
	% https://q.uiver.app/#q=WzAsNSxbMCwwLCJUKFgpIl0sWzEsMCwiVChYJykiXSxbMCwxLCJYIl0sWzEsMSwiWCciXSxbMiwxLCJRIl0sWzAsMiwiXFx4aSIsMl0sWzEsMywiXFx4aSciXSxbMCwxLCJUKHUpIiwwLHsib2Zmc2V0IjotMX1dLFswLDEsIlQodikiLDIseyJvZmZzZXQiOjF9XSxbMiwzLCJ1IiwwLHsib2Zmc2V0IjotMX1dLFsyLDMsInYiLDIseyJvZmZzZXQiOjF9XSxbMyw0LCJxIiwyXV0=
	\[\begin{tikzcd}
		{T(X)} & {T(X')} & \\
		X & {X'} & Q
		\arrow["{T(u)}", shift left, from=1-1, to=1-2]
		\arrow["{T(v)}"', shift right, from=1-1, to=1-2]
		\arrow["\xi"', from=1-1, to=2-1]
		\arrow["{\xi'}", from=1-2, to=2-2]
		\arrow["u", shift left, from=2-1, to=2-2]
		\arrow["v"', shift right, from=2-1, to=2-2]
		\arrow["q"', from=2-2, to=2-3]
	\end{tikzcd}\]	
	Since $ q$ is an absolute coequaliser of $ u$ and $ v$ in $ \mathbb{X}$, then $ T(q)$ is a coequaliser of $ T(u)$ and $ T(v)$ in $ \mathbb{A}$. 
	Since the $u$- and $ v$-diagrams commute, then $ q\xi'$ also coequalises $ T(u)$ and $ T(v)$. 
	Then there is a unique morphism $k: T(Q) \longrightarrow Q$ such that $ kT(q) = q \xi'$.
	Since this square commutes, then $ q$ is also a morphism in $ \mathbb{X}^{T}$.
	Now we want to show that $ q$ is the coequaliser of $ u$ and $ v$ in $ \mathbb{X}^{ T}$.
	Let $ q:(X',\xi') \longrightarrow (Q', k')$ be a morphism of $ T$-algebras such that $ q'u = q'v$ in $ \mathbb{X}^{T}$.
	Then $ q'u=q'v $ in $ \mathbb{X}$, and $ T(q')T(u)=T(q')T(v)$ in $ \mathbb{X}$, which, respectively, induce $l$ and $ T(l)$ in the diagram below, by universal property of the coequalisers.
	
	% https://q.uiver.app/#q=WzAsOCxbMCwwLCJUKFgpIl0sWzIsMCwiVChYJykiXSxbMCwyLCJYIl0sWzIsMiwiWCciXSxbMywyLCJRIl0sWzMsMCwiVChRKSJdLFs0LDEsIlQoUScpIl0sWzQsMywiUSciXSxbMCwyLCJcXHhpIiwyXSxbMSwzLCJcXHhpJyJdLFswLDEsIlQodSkiLDAseyJvZmZzZXQiOi0xfV0sWzAsMSwiVCh2KSIsMix7Im9mZnNldCI6MX1dLFsyLDMsInUiLDAseyJvZmZzZXQiOi0xfV0sWzIsMywidiIsMix7Im9mZnNldCI6MX1dLFszLDQsInEiXSxbMSw1LCJUKHEpIl0sWzEsNiwiVChxJykiLDIseyJsYWJlbF9wb3NpdGlvbiI6MzB9XSxbMyw3LCJxJyIsMl0sWzUsNCwiayIsMCx7InN0eWxlIjp7ImJvZHkiOnsibmFtZSI6ImRhc2hlZCJ9fX1dLFs2LDcsImsnIiwwLHsic3R5bGUiOnsiYm9keSI6eyJuYW1lIjoiZGFzaGVkIn19fV0sWzUsNiwiVChsKSJdLFs0LDcsImwiXV0=
	\[\begin{tikzcd}
		{T(X)} && {T(X')} & {T(Q)} & \\
		&&&& {T(Q')} \\
		X && {X'} & Q \\
		&&&& {Q'}
		\arrow["{T(u)}", shift left, from=1-1, to=1-3]
		\arrow["{T(v)}"', shift right, from=1-1, to=1-3]
		\arrow["\xi"', from=1-1, to=3-1]
		\arrow["{T(q)}", from=1-3, to=1-4]
		\arrow["{T(q')}"'{pos=0.3}, from=1-3, to=2-5]
		\arrow["{\xi'}", from=1-3, to=3-3]
		\arrow["{T(l)}", dashed, from=1-4, to=2-5]
		\arrow["k", dashed, from=1-4, to=3-4]
		\arrow["{k'}", from=2-5, to=4-5]
		\arrow["u", shift left, from=3-1, to=3-3]
		\arrow["v"', shift right, from=3-1, to=3-3]
		\arrow["q", from=3-3, to=3-4]
		\arrow["{q'}"', from=3-3, to=4-5]
		\arrow["l", dashed, from=3-4, to=4-5]
	\end{tikzcd}\]
	Now we just need to show that $l$ is also a morphism of $T$-algebras, which is to show that $k'T(l)=lk$. 
	Consider a segment of the diagram:
	% https://q.uiver.app/#q=WzAsNixbMCwwLCJUKFgnKSJdLFswLDEsIlgnIl0sWzEsMCwiVChRKSJdLFsxLDEsIlEiXSxbMiwwLCJUKFEnKSJdLFsyLDEsIlEnIl0sWzAsMiwiVChxKSJdLFsxLDMsInEiLDJdLFswLDEsIlxceGknIiwyXSxbMiwzLCJrIl0sWzIsNCwiVChsKSIsMCx7InN0eWxlIjp7ImJvZHkiOnsibmFtZSI6ImRhc2hlZCJ9fX1dLFszLDUsImwiLDIseyJzdHlsZSI6eyJib2R5Ijp7Im5hbWUiOiJkYXNoZWQifX19XSxbNCw1LCJrJyJdXQ==
	\[\begin{tikzcd}
		{T(X')} & {T(Q)} & {T(Q')} \\
		{X'} & Q & {Q'}
		\arrow["{T(q)}", from=1-1, to=1-2]
		\arrow["{\xi'}"', from=1-1, to=2-1]
		\arrow["{T(l)}", dashed, from=1-2, to=1-3]
		\arrow["k", from=1-2, to=2-2]
		\arrow["{k'}", from=1-3, to=2-3]
		\arrow["q"', from=2-1, to=2-2]
		\arrow["l"', dashed, from=2-2, to=2-3]
	\end{tikzcd}\]
	
	The outer square commutes because $ k' T(l)T(q)=k' T(q')=q'\xi'= lq \xi'$. 
	The left square also commutes by assumption, and T(q) is an epi since it is a coequaliser.
	Then by a pasting lemma, the right square commutes, too.
	Then $ l$ is indeed a $T$-algebra morphism, which shows that $ q$ is the coequaliser of $ u,v$ in $ \mathbb{X}^{T}$.
	From $ (1)$, $ \mathbb{X}^{T}\sim \mathbb{A}$, which means there is also a coequaliser of the diagram $ \mathcal{D}$ in $ \mathbb{A}$, which is trivially preserved by $ U$.
\end{proof}


Now we want a simpler way to check monadicity, which will become useful in the next example.

\begin{thm}\label{thm:crude}{ \textit{Crude Monadicity Theorem}}
	Let $ U: \mathbb{A} \longrightarrow \mathbb{X}$ be a functor. $ U$ is monadic if
	\begin{enumerate}
		\item $ U$ has a left adjoint;
		\item $ \mathbb{A}$ has, and $ U$ preserves, reflexive coequalisers;
		\item $ U$ reflects isomorphisms.
	\end{enumerate}
	
\end{thm}

% TODO Finish this proof in full
\begin{proof}
	Let us only prove the second point. 
\end{proof}

